\begin{table*}[t]
\centering\small
\begin{tabular}{rrrrr}
\toprule
$q$ & $\alpha$ & Risk & Sel. (\%) & Over / 40 \\
\midrule
0.05 & 0.02 & 0.003 & 48.3 & 0 \\
0.05 & 0.05 & 0.033 & 66.1 & 3 \\
0.05 & 0.10 & 0.078 & 96.4 & 1 \\
0.05 & 0.15 & 0.083 & 100.0 & 0 \\
0.05 & 0.20 & 0.083 & 100.0 & 0 \\
0.10 & 0.02 & 0.011 & 57.8 & 0 \\
0.10 & 0.05 & 0.040 & 77.6 & 4 \\
0.10 & 0.10 & 0.070 & 100.0 & 0 \\
0.10 & 0.15 & 0.070 & 100.0 & 0 \\
0.10 & 0.20 & 0.070 & 100.0 & 0 \\
0.20 & 0.02 & 0.016 & 71.9 & 4 \\
0.20 & 0.05 & 0.045 & 95.8 & 3 \\
0.20 & 0.10 & 0.050 & 100.0 & 0 \\
0.20 & 0.15 & 0.050 & 100.0 & 0 \\
0.20 & 0.20 & 0.050 & 100.0 & 0 \\
0.30 & 0.02 & 0.017 & 80.2 & 8 \\
0.30 & 0.05 & 0.041 & 100.0 & 0 \\
0.30 & 0.10 & 0.041 & 100.0 & 0 \\
0.30 & 0.15 & 0.041 & 100.0 & 0 \\
0.30 & 0.20 & 0.041 & 100.0 & 0 \\
\bottomrule
\end{tabular}
\caption{RewardBench 2, median pooling: grouped SCoRE-MDR for the same churn loss and target $\alpha/B$. Risk and selected coverage $K/|S_0|$ average 40 splits. Over/40 counts splits whose realized churn exceeds the budget. Calibration candidate scores are fully observed; selection uses no candidate test scores. Guarantee assumptions and score-entry accounting appear in Appendix~\ref{app:review-score} and Table~\ref{tab:review-workload}.}
\label{tab:review-score-0-median}
\end{table*}
